\documentclass[a4paper,11pt]{article}
\usepackage[utf8]{inputenc}
\usepackage[T1]{fontenc}
\usepackage[french]{babel}
\usepackage{amsmath,amsfonts,amssymb}
\usepackage{geometry}
\usepackage{fancyhdr}
\usepackage{titlesec}
\usepackage{enumitem}
\usepackage{listings}

\geometry{margin=2.5cm}

% Header and footer configuration
\pagestyle{fancy}
\fancyhf{}
\fancyhead[L]{Licence 1 Informatique}
\fancyhead[R]{CERI - Avignon Université}
\fancyfoot[C]{Mathématiques Discrètes \hfill TD3 - Logique de premier ordre \hfill Page \thepage/3}

% Title formatting
\titleformat{\section}
{\Large\bfseries}
{\thesection}
{1em}
{}

\titleformat{\subsection}
{\large\bfseries}
{\thesubsection}
{1em}
{}

% Configuration listings pour pseudo-code
\lstset{
    basicstyle=\ttfamily\small,
    keywordstyle=\bfseries,
    commentstyle=\itshape\color{gray},
    frame=single,
    breaklines=true,
    showstringspaces=false
}

\begin{document}

\begin{center}
{\Huge\bfseries TD3 - Logique de premier ordre} \\[0.5cm]
{\large Hiver 2026} \\[0.3cm]
\end{center}

\vspace{0.5cm}

L'objectif de ce TD est de mettre en pratique les bases de la logique de premier ordre.

\section{De la logique de premier ordre au français}

Traduisez en français les formules logiques suivantes où les prédicats signifient :

$E(x)$ : ``$x$ est étudiant.'' \quad $P(x)$ : ``$x$ a visité Paris.'' \quad $C(x)$ : ``$x$ a un chat.'' \quad $F(x)$ : ``$x$ a un forêt.''

\begin{enumerate}
\item $\forall x(P(x) \wedge C(x))$
\item $\exists x(C(x) \wedge \neg F(x))$
\item $\forall x(E(x) \wedge P(x) \Rightarrow C(x))$
\item $\forall x(P(x) \Rightarrow F(x))$
\item $\neg \exists x(E(x) \wedge C(x) \wedge F(x))$
\end{enumerate}

\section{Du français à la logique de premier ordre}

\subsection{Traduire des phrases}

Donnez la formule logique correspondant aux phrases suivantes en utilisant le prédicat suivant :

$D(x,y)$ : ``La personne $x$ peut duper la personne $y$''

On considérera que le domaine des variables est \emph{tout le monde}.

\begin{enumerate}
\item Tout le monde peut duper Alban.
\item Brigitte peut duper tout le monde.
\item Tout le monde peut duper quelqu'un.
\item Il n'y a personne qui puisse duper tout le monde.
\item Tout le monde peut être dupé par quelqu'un.
\item Personne ne peut duper à la fois Alban et Brigitte.
\item Personne ne peut se duper lui-même.
\item Il y des gens qui peuvent duper aucune personne à part eux-mêmes.
\end{enumerate}

\subsection{Quelques négations}

9. Pour chacune des phrases précédentes, donnez la négation de la formule obtenue, puis donnez la phrase en français correspondant à cette négation.

\section{Identifier des inférences}

Pour chacun des cas, utilisez un formalisme logique pour traduire les hypothèses et la conclusion et précisez la règle d'inférence utilisée pour obtenir cette conclusion à partir des hypothèses.

\begin{enumerate}
\item \emph{``Tous les hommes sont mortels. Or Socrate est un homme.''}\\
On en déduit : \emph{``Socrate est mortel.''}

\item \emph{``Alice aime les mathématiques.''}\\
On en déduit : \emph{``Il y a des gens qui aiment les mathématiques.''}

\item \emph{``Tous les gens qui n'aiment pas les burgers aiment les haricots. Tous les gens qui aiment les haricots aiment le sport.''}\\
On en déduit : \emph{``Tous les gens qui n'aiment pas les burgers aiment le sport.''}

\item \emph{``Tout le monde aime les clafoutis.''}\\
On en déduit : \emph{``Wrek aime le clafoutis.''}

\item \emph{``Tous les gens qui ont un ordinateur savent envoyer un mail. Mon grand-père n'a pas d'ordinateur.''}\\
On en déduit : \emph{``Mon grand-père ne sait pas envoyer de mails.''}
\end{enumerate}

\section{Modéliser et déduire}

Pour chacun des cas, utilisez un formalisme logique pour traduire les hypothèses données et la conclusion, puis utilisez les règles d'inférence pour procéder à une argumentation afin de conclure.

\begin{enumerate}
\item \emph{``Tous les colibris ont des couleurs vives. Aucun gros oiseau ne se nourrit de miel. Les oiseaux qui ne se nourrissent pas de miel ont des couleurs ternes.''}

Montrez qu'on peut en déduire : \emph{``Les colibris sont des petits oiseaux.''}

\item \emph{``Tous les gens qui font du sport mangent sainement et boivent beaucoup d'eau. Tous les étudiants de l'université font du sport et dorment bien la nuit et boivent beaucoup d'eau.''}

Montrez qu'on peut en déduire : \emph{``Tous les étudiants de l'université dorment bien la nuit et boivent beaucoup d'eau.''}

\item \emph{``Tout le monde a un chat ou un oiseau. Tous les gens qui ont un oiseau mais pas de chat ont un poisson rouge.''}

Montrez qu'on peut en déduire : \emph{``Tous les gens qui n'ont pas de poisson rouge ont un chat.''}

\item \emph{``Tout le monde aime voyager ou rester chez soi et écouter des musiques. Tout le monde aime faire du sport ou ne pas écouter de musique. Je connais une personne qui n'aime pas voyager.''}

Montrez qu'on peut en déduire : \emph{``Il y a des gens qui n'aiment pas le chant.''}

\item \emph{``Au zoo, Koko est un singe. Tous les animaux du zoo aiment jouer ou aiment les légumes. Tous les animaux qui aiment les légumes aiment les légumes. Les singes détestent le poulet.''}

Montrez qu'on peut en déduire : \emph{``Il y a des animaux qui aiment jouer dans ce zoo.''}
\end{enumerate}

\section{Analyser des raisonnements}

Pour chacun des cas, précisez si le raisonnement est correct en justifiant votre réponse.

\begin{enumerate}
\item ``Tous les étudiants de première année suivent un cours de programmation. Alice est étudiante de première année.''

Peut-on en déduire qu'Alice suit des cours de programmation ?

\item ``Tous les étudiants d'informatique suivent un cours de programmation. Bob suit des cours de programmation.''

Peut-on en déduire que Bob est étudiant d'informatique ?

\item ``Tous les perroquets aiment les fruits. Tui n'est pas un perroquet.''

Peut-on en déduire que Tui n'aime pas les fruits ?

\item ``Toutes celles et ceux qui mangent une pomme chaque jour sont en bonne santé. Carole n'est pas en bonne santé.''

Peut-on en déduire que Carole ne mange pas de pomme chaque jour ?

\item ``Dan aime tous les films d'action. Dan aime le film 'Die Hard'.''

Peut-on en déduire que 'Die Hard' est un film d'action ?
\end{enumerate}

\section{Montrer des règles d'inférences}

Nous supposons ici que l'on connaît et peut appliquer les règles d'inférence de la logique propositionnelle ainsi que les règles d'instanciation universelle, de généralisation universelle, d'instanciation existentielle et de généralisation existentielle.

\begin{enumerate}
\item Montrez que la règle de modus ponens universel est correcte.
\item Montrez que la règle de modus tollens universel est correcte.
\item Montrez que la règle de transitivité universelle est correcte.
\end{enumerate}

\section{Algorithmiques et quantificateurs}

On donne le pseudo-code de deux procédures p1 et p2 :

\begin{minipage}[t]{0.48\textwidth}
\textbf{Procedure p1()}
\begin{lstlisting}
bool <- vrai
PourTout i ∈ A Faire
  PourTout j ∈ B Faire
    bool <- bool∧F(i,j)
  FinPourTout
FinPourTout
Retourner bool
FinProcedure
\end{lstlisting}
\end{minipage}
\hfill
\begin{minipage}[t]{0.48\textwidth}
\textbf{Procedure p2()}
\begin{lstlisting}
bool <- vrai
PourTout i ∈ A Faire
  bool2 <- faux
  PourTout j ∈ B Faire
    bool2 <- bool2∧P(i,j)
  FinPourTout
  bool <- bool∧bool2
FinPourTout
Retourner bool
FinProcedure
\end{lstlisting}
\end{minipage}

\vspace{0.5cm}

\begin{enumerate}
\item Observez le pseudo-code de la procédure p1. Donnez la formule logique qui est vraie si et seulement si p1 renvoie \texttt{vrai}.

\item Observez le pseudo-code de la procédure p2. Donnez la formule logique qui est vraie si et seulement si p2 renvoie \texttt{vrai}.

\item Donnez le pseudo-code de la procédure p3, qui renvoie \texttt{vrai} si et seulement si la formule $\exists x \forall y P(x,y)$ est vraie. On considérera que le domaine de la variable $x$ est l'ensemble $A$ et que le domaine de la variable $y$ est l'ensemble $B$.

\item Donnez le pseudo-code de la procédure p4, qui renvoie \texttt{vrai} si et seulement si la formule $\exists x \exists y P(x,y)$ est vraie. On considérera que le domaine de la variable $x$ est l'ensemble $A$ et que le domaine de la variable $y$ est l'ensemble $B$.
\end{enumerate}

\end{document}
